\section{Answer 3: the definition of a path between \texorpdfstring{$n_i$}{n_i} and \texorpdfstring{$n_j$}{n_j}}\label{sec:answer3}
% claims used here: C12 C15 C16 C17 C21 C22 C23 C24 C25 m2-5 m2-7 m2-8 m2-13 m2-15

Question~3 asks for the definition of a path between two vertices $n_i$ and
$n_j$. The answer is \cref{def:path}, restated for a pair of vertices in
\cref{def:pathbetween} below: in the course, a path is a walk with no repeated vertex,
which is the textbook's Definition~1.2.13~\cite[Definition~1.2.13, p.~22]{petrosyan2015grafneri}.
This section adds only the form for a pair of vertices and the comparison of the three
naming conventions. The board uses the word in its definition of ``connected'' and does not
say which meaning it has~\cite{whiteboard1,whiteboard2}. The assignment text is not reproduced in this paper. As reported by the user, it states that in the course ``path'' means a walk with no repeated vertices; the answer to Question~3 rests on this statement, so it is \kindtag{interpretation, qualified}: it is the user's report, and the board does not confirm it.

\begin{plainwords}
A path between two houses is a route from one to the other. It never passes a
house twice, so it has no detour. If the two houses are the same, the route has no
roads at all. It is just the one house.
\end{plainwords}

\begin{definition}[Path between $n_i$ and $n_j$, course convention]\label{def:pathbetween}
Let $n_i$ and $n_j$ be vertices of $G$. A \emph{path between $n_i$ and $n_j$} is a path in the sense of \cref{def:path} whose first vertex is $n_i$ and whose last vertex is $n_j$: a walk $w_0,w_1,\dots,w_k$ with $w_0=n_i$, $w_k=n_j$ and pairwise distinct vertices. Its length is $k$.
\end{definition}

This definition is taken from the assignment text, as reported by the user, together with the textbook's Definition~1.2.13; it is not derived.

This is a simple path in the textbook's sense and a path in English usage~\cite{petrosyan2015grafneri,grinberg2025graph}; Grinberg phrases the pair form as ``a path from $p$ to $q$ is a path that starts at $p$ and ends at $q$''~\cite[Definition~2.9.1(d), p.~36]{grinberg2025graph}.

\begin{remark}\label{rem:pathbetween}
The definition forbids repeated vertices, and so repeated edges (\cref{prop:pathtrail}). For $i=j$ the only such path is the single vertex $n_i$, of length $0$ (this uses the reading of length-$0$ walks in \cref{sec:conventions}, item~7). For $i\ne j$ the length is at least $1$ and, by \cref{cor:bound}, at most $N-1$. In the board's sentence ``there exists a path between nodes'', the words mean that at least one such object exists. By \cref{prop:readings} that statement has the same truth value whether ``path'' means walk, trail or path in the sense above. What changes with the convention is how many such objects there are (\cref{rem:howmany}).
\end{remark}

\begin{example}\label{ex:pathbetween}
\Cref{ex:path} gives the path $1,2,3$ between $n_1$ and $n_3$ in $B$. It is the only one: every path has at most $N-1=2$ edges (\cref{cor:bound}), and the enumeration of walks up to length $4$ finds no other (by script). Between $n_1$ and $n_1$ the only path is the single vertex $1$.
\end{example}

\begin{nonexample}\label{nex:pathbetween}
The walk $1,2,1,2,3$ of \cref{ex:walk} joins $n_1$ and $n_3$ but is not a path between them (\cref{nex:path}, \cref{nex:trail}). In the graph $D$ there is no path between $n_1$ and $n_3$ at all.
\end{nonexample}

\subsection{The three conventions in one table}\label{sec:conventionstable}

\begin{plainwords}
The same two rules get different names in different books. One rule says never
drive a road twice. The other says never pass a house twice. Our textbook, English
books and our course do not always use the same word for the same rule. So here
is one table that puts the words side by side.
\end{plainwords}

\begin{table}[ht]
\centering
\caption{The three naming conventions. ``Course (stated)'' is what the assignment text states, as reported by the user: only the word path. For every other row the course convention is not stated, so this paper follows the textbook and writes ``simple cycle'' for the English ``cycle''. Sources: textbook~\cite{petrosyan2015grafneri}, Grinberg~\cite{grinberg2025graph}, Wikipedia~\cite{wikipedia_path_graph_theory}.}
\label{tab:conventions}
\small
\begin{tabular}{@{}p{3.6cm}p{2.8cm}p{2.4cm}p{2.2cm}p{2.6cm}@{}}
\toprule
Property of the walk & Textbook word (number) & English word & Course (stated) & This paper \\
\midrule
any list of steps along edges & walk (1.2.10) & walk & not stated & walk \\
walk with first vertex equal to last & closed walk (1.2.11) & closed walk & not stated & closed walk \\
no repeated edge; vertices may repeat & ``path'' (1.2.12) & trail & not stated & \textbf{trail} \\
no repeated vertex (so no repeated edge) & ``simple path'' (1.2.13) & path & path & \textbf{path} \\
closed, no repeated edge; vertices may repeat & ``cycle'' or ``closed path'' (1.2.14) & closed trail & not stated & cycle (textbook meaning) \\
closed, no repeated vertex except first and last, length at least $3$ & ``simple cycle'' (1.2.15) & cycle & not stated & simple cycle \\
\bottomrule
\end{tabular}
\end{table}

The textbook's rows are its Definitions~1.2.10 to~1.2.15~\cite{petrosyan2015grafneri}. The English words for walk, trail, path and cycle are Grinberg's Definitions~2.9.1(c), 3.4.5 and 2.10.1(b)~\cite{grinberg2025graph}; ``closed trail'' is this paper's own name for the English counterpart of the textbook's cycle, since the sources checked give no single standard word for it. Wikipedia says that a path has distinct vertices ``by most definitions''~\cite{wikipedia_path_graph_theory}, and also records that some authors use ``path'' for walks with repeated vertices and ``simple path'' for those with distinct vertices~\cite{wikipedia_path_graph_theory}. That is the textbook's pattern, so the course convention is one of at least two in use.

\begin{plainwords}
The word changes how many routes we count. It does not change whether a route
exists. In the board village there is one route from the first house to the
last. But there are three trips of at most four roads, and two of them turn back.
\end{plainwords}

\begin{remark}[The number of objects depends on the convention]\label{rem:howmany}
Under the course convention, the number of paths from $n_1$ to $n_3$ in $B$ is $1$. The walks of length at most $4$ from $n_1$ to $n_3$ are $1,2,3$ (length $2$), $1,2,1,2,3$ and $1,2,3,2,3$ (length $4$); only the first is a path or a trail (by script). So the number of walks, the number of trails and the number of paths are different numbers in general, and the entries of $A^k$ give only the first (\cref{prop:warning}). The existence statement used by \cref{def:connected} is the same under every convention (\cref{prop:readings}).
\end{remark}
